\documentclass[,a4paper]{article}

% --- Pacotes Necessários ---
\usepackage[utf8]{inputenc}
\usepackage[T1]{fontenc}
\usepackage[portuguese]{babel}
\usepackage[margin=0cm]{geometry}
\usepackage{xcolor}
\usepackage{fontawesome5}
\usepackage{tikz}
\usepackage{graphicx}
\usepackage{enumitem}
\usepackage{multicol}
\usepackage{hyperref}
\usepackage{accsupp}
\hypersetup{hidelinks}
\usetikzlibrary{shapes, shadows, arrows, positioning}

% --- Configuração de Grid Compacto ---
\setlength{\multicolsep}{2pt}
\setlength{\columnsep}{5pt} % <--- REDUZ O ESPAÇO HORIZONTAL ENTRE COLUNAS

% --- Definição de Cores ---
\definecolor{mainpurple}{HTML}{8B3DEC}
\definecolor{sidebargray}{HTML}{1F1F1F}
\definecolor{lightbg}{HTML}{F8F9FA}
\definecolor{textgray}{HTML}{444444}
\definecolor{darkpurple}{HTML}{6A2EB5}

% --- Configurações de Fonte ---
\usepackage[sfdefault]{roboto}
\hbadness=10000

% --- Título da Sidebar (alinhamento visual) ---
\newcommand{\sidebartitle}[2]{%
  \begin{center}
    {\color{mainpurple}\bfseries\LARGE \makebox[1.3em][c]{#1}\hspace{0.35em}#2}
  \end{center}%
  \vspace{0.14cm}%
}
\newcommand{\sideblockspace}{\vspace{0.1cm}}

\begin{document}

\pagestyle{empty}

% --- Estrutura de Fundo (TikZ) ---
\begin{tikzpicture}[remember picture, overlay]
  \fill[sidebargray] (current page.north west) rectangle ([xshift=7.5cm]current page.south west);
  \fill[mainpurple] ([xshift=7cm, yshift=-1cm]current page.north west)
  [rounded corners=20pt] -- ([xshift=20.5cm, yshift=-1cm]current page.north west)
  -- ([xshift=20.5cm, yshift=-4.2cm]current page.north west)
  [rounded corners=20pt] -- ([xshift=7.5cm, yshift=-4.2cm]current page.north west) -- cycle;
  \node[circle, draw=white, line width=3pt, minimum size=3.5cm, inner sep=0pt, fill=lightbg,
    path picture={
      \IfFileExists{assets/profile.png}{
        \begin{scope}
          \clip (path picture bounding box.center) circle [radius=1.75cm];
          \node at (path picture bounding box.center) {
            \includegraphics[width=3.5cm,height=3.5cm]{assets/profile.png}
          };
        \end{scope}
      }{
        \node at (path picture bounding box.center) {\color{textgray}\footnotesize FOTO};
      }
    }
  ] at ([xshift=7.5cm, yshift=-4.2cm]current page.north west) {};
\end{tikzpicture}

% --- Conteúdo ---
\begin{minipage}[t]{6.5cm} % COLUNA LATERAL
  \vspace{1.2cm}
  \begin{flushleft}
    \color{white}
    \sidebartitle{\faUser}{CONTATO}
    \vspace{0.6cm}
    \faLinkedin \quad \href{https://linkedin.com/in/renanrod4}{renanrod4} \\
    \vspace{0.2cm}
    \faPhone* \quad (11) 93340-7053 \\
    \vspace{0.2cm}
    \faGithub \quad \href{https://github.com/renanrod4}{github/renanrod4} \\
    \vspace{0.2cm}
    \faGlobe \quad \href{https://renanrod.vercel.app/}{renanrod.vercel.app} \\
    \vspace{0.2cm}

    \faEnvelope\quad \href{mailto:renanrdemeneses@gmail.com}{\normalsize renanrdemeneses@gmail.com} \\
    \vspace{0.2cm}
    \faMapMarker* \quad São Paulo, SP - Brasil\\

    \vspace{0.5cm}

    \sidebartitle{\faCode}{SKILLS}
    \vspace{-0.3cm}

    \newcommand{\skilltag}[1]{
      \tikz[baseline]\node[
        fill=mainpurple!30!black,
        draw=mainpurple!75,
        rounded corners=4pt,
        inner xsep=7pt,
        inner ysep=4pt,
        text=lightbg,
        font=\normalsize\bfseries
      ] {#1};
    }

    \newcommand{\langtag}[3]{%
      \tikz[baseline]\node[
        fill=mainpurple!30!black,
        draw=mainpurple!75,
        rounded corners=4pt,
        inner xsep=8pt,
        inner ysep=4.5pt,
        text=lightbg,
        font=\normalsize\bfseries
      ] {\raisebox{-0.2em}{\includegraphics[height=1.25em]{#1}}\hspace{0.45em}\strut #2\hspace{0.4em}{\small\strut #3}};%
    }

    \begin{center}
      \normalsize\textbf{Software \& Ferramentas} \\
      \vspace{0.08cm}
      \skilltag{Blender}\hspace{2pt}\skilltag{Unity}\hspace{2pt}\skilltag{VS Code}\\
      \vspace{0.08cm}
      \skilltag{Arduino}\hspace{2pt}\skilltag{Figma}\hspace{2pt}\skilltag{Git}\\
      \vspace{0.08cm}
      \skilltag{GitHub}\hspace{2pt}\skilltag{Bootstrap}\hspace{2pt}\skilltag{Tailwind}\\
      \vspace{0.08cm}
      \skilltag{Vercel}
    \end{center}

    \sideblockspace
    \begin{center}
      \normalsize\textbf{Desenvolvimento Web} \\
      \vspace{0.08cm}
      \skilltag{HTML}\hspace{2pt}\skilltag{CSS}\hspace{2pt}\skilltag{JavaScript}\\
      \vspace{0.08cm}
      \skilltag{TypeScript}\hspace{2pt}\skilltag{React}\hspace{2pt}\skilltag{Next.js}\\
      \vspace{0.08cm}
      \skilltag{Node.js}\hspace{2pt}\skilltag{PHP}\hspace{2pt}\skilltag{REST}
    \end{center}

    \sideblockspace
    \begin{center}
      \normalsize\textbf{Linguagens} \\
      \vspace{0.08cm}
      \skilltag{Python}\hspace{2pt}\skilltag{Rust}\hspace{2pt}\skilltag{C++}\\
      \vspace{0.08cm}
      \skilltag{C\#}\hspace{2pt}\skilltag{Bash}
    \end{center}

    \sideblockspace
    \begin{center}
      \normalsize\textbf{Bancos de Dados} \\
      \vspace{0.08cm}
      \skilltag{MySQL}\hspace{2pt}\skilltag{MongoDB}
    \end{center}

    \vspace{0.3cm}

    \sidebartitle{\faGlobe}{IDIOMAS}
    \vspace{0cm}
    \begin{center}
      \begin{tabular}{@{}c@{}}
        \langtag{assets/flags/br.png}{Português}{Nativo} \\
        [0.6cm]
        \langtag{assets/flags/us.png}{Inglês}{Avançado} \\
        [0.6cm]
        \langtag{assets/flags/de.png}{Alemão}{Básico}
      \end{tabular}
    \end{center}
  \end{flushleft}
\end{minipage}
\hspace{1.5cm}
\begin{minipage}[t]{11.5cm} % COLUNA PRINCIPAL
  \vspace{1.2cm}
  {\Huge \textbf{\color{white} RENAN RODRIGUES}} \\
  \vspace{0.2cm}
  {\hspace*{1cm}\large\textbf{ \color{white!80} Engenharia da Computação | Full Stack Developer}}

  \vspace{1.4cm}

  \begin{minipage}{11.0cm}
    \section*{\hspace*{0.5cm}\vspace{-0.4cm}\color{mainpurple}{\faIdCard \quad RESUMO}}
    \vspace{0.2cm}
    {\normalsize \color{textgray} Estudante de Engenharia da Computação focado em desenvolvimento Full Stack. Experiência na criação de aplicações web, APIs REST e sistemas interativos utilizando JavaScript/TypeScript, React, Next.js e Rust. Busco aplicar meus conhecimentos em projetos reais e evoluir continuamente como desenvolvedor.\par}
    \vspace*{-0.2cm}
    \section*{\hspace*{0.5cm}\color{mainpurple}{\faGraduationCap \quad FORMAÇÃO}}
    \noindent \textbf{Bacharelado em Engenharia de Computação} \hfill \textbf{2025 -- 2030} \\
    \textit{\color{textgray}Universidade de Sorocaba (UNISO)}\par
    \vspace{0.3cm}
    \noindent \textbf{Ensino Médio Completo} \hfill \textbf{2022 -- 2024} \\
    \textit{\color{textgray}E.E. Bel Elias Alves da Costa Bacharel}{\color{textgray}\hfill{2023 -- 2024}}\\
    \textit{\color{textgray}Colégio Cognos Sistema Mackenzie de Ensino} {\color{textgray}\hfill{2022}}
    \vspace*{-0.2cm}
    \section*{\hspace*{0.5cm}\color{mainpurple}{\faFlask \quad EXPERIÊNCIA ACADÊMICA E P\&D (UNISO)}}
    {\normalsize
    \noindent {\color{black}\textbf{Uniso Flow} | \textit{Full Stack Developer} \hfill \textbf{2025 -- Presente}}\\
    \color{textgray}Análise preditiva para fluxo acadêmico, com cruzamento de pré-requisitos para prever demanda e gargalos.\\[-0.05cm]
    \color{textgray}\textbf{Tech:} Next.js, Tailwind, Radix UI, Tauri (Rust).\\[0.1cm]
    \noindent {\color{black}\textbf{Termômetro Laser Industrial} | \textit{ESP32 \& C++} \hfill \textbf{2025 -- Presente}}\\
    \color{textgray}Protótipo para medição térmica de precisão e firmware para calibração óptica/processamento de sinais.\\[-0.05cm]
    \color{textgray}\textbf{Tech:} C++, ESP32, Sensores Infravermelhos.
    }
    \vspace*{-0.2cm}
    \section*{\hspace*{0.5cm}\color{mainpurple}{\faRocket \quad PROJETOS TÉCNICOS}}

    \href{https://octodev.vercel.app/}{\tikz[baseline]\node[fill=mainpurple!10, rounded corners=8pt, inner sep=6pt] {
        \begin{minipage}{10.55cm}
          \raggedright
          \textbf{\color{darkpurple}OctoDev} | \textit{React, Next.js, MongoDB, IA}\\
          \small Plataforma educacional Full Stack para ensino de programação gamificado, com desafios, múltiplas linguagens e correção de código em tempo real.
        \end{minipage}
    };}
    \vspace{0.15cm}

    \href{https://helio-sync.vercel.app}{\tikz[baseline]\node[fill=mainpurple!10, rounded corners=8pt, inner sep=6pt] {
        \begin{minipage}{10.55cm}
          \raggedright
          \textbf{\color{darkpurple}Helio Sync} | \textit{TypeScript, Python, MQTT}\\
          \small Plataforma mecatrônica para rastreamento solar em dois eixos, com telemetria em tempo real, análise de eficiência e controle remoto do hardware.
        \end{minipage}
    };}
    \vspace{0.15cm}

    \href{https://github.com/renanrod4/vscode-code-in-comments}{\tikz[baseline]\node[fill=mainpurple!10, rounded corners=8pt, inner sep=6pt] {
        \begin{minipage}{10.55cm}
          \raggedright
          \textbf{\color{darkpurple}CodeInComments} | \textit{VS Code Extension, JavaScript, Rust, Python}\\
          \small Extensão do VS Code que destaca blocos de código dentro de comentários, preservando o realce de sintaxe e oferecendo efeitos configuráveis.
        \end{minipage}
    };}
    \vspace*{-0.2cm}
    \section*{\hspace*{0.5cm}\color{mainpurple}{\faTrophy \quad CONQUISTAS}}
    \begin{itemize}[label=\color{mainpurple}\faCheck, leftmargin=1.2em, itemsep=4pt]
      \item \textbf{Feira de Ciências de SP (2023):} Pesquisa de múons (Fermilab/USP).
      \item \textbf{Medalha de Bronze:} Olimpíada de Matemática (2024).
      \item \textbf{Medalhista:} Olimpíada Brasileira de Astronomia (2024).
    \end{itemize}
  \end{minipage}

\end{minipage}

% --- Footer fixo no rodapé da página ---
\begin{tikzpicture}[remember picture, overlay]
  \node[anchor=south, yshift=0.35cm] at ([xshift=14.25cm]current page.south west) {
    \footnotesize \color{textgray} \faCode \hspace{2pt} Desenvolvido em \LaTeX \hspace{1pt} por Renan Rodrigues
  };
\end{tikzpicture}

\end{document}
